\documentclass[11pt,a4paper]{article}
\usepackage[utf8]{inputenc}
\usepackage[T1]{fontenc}
\usepackage[catalan]{babel}
\usepackage{lmodern}
\usepackage[sfdefault,lining]{FiraSans}
\usepackage[a4paper,top=2.3cm,bottom=2.2cm,left=1.8cm,right=1.8cm,headheight=1.3cm,headsep=0.4cm,footskip=1cm]{geometry}
\usepackage{amsmath,amssymb}
\usepackage{icomma}
\usepackage{graphicx}
\usepackage[table]{xcolor}
\usepackage{tikz}
\usepackage{fontawesome5}
\usepackage{ccicons}
\usepackage{fancyhdr}
\usepackage{lastpage}
\usepackage{tasks}
\usepackage{needspace}
\usepackage[most]{tcolorbox}
\usepackage{array}
\usepackage[hidelinks,pdftitle={Solucions de la fitxa de la unitat 8: Distribucions bidimensionals. 1r de Batxillerat CCSS},pdfsubject={Matemàtiques Aplicades a les CCSS I},pdfkeywords={distribucions bidimensionals, taula de doble entrada, núvol de punts, correlació, regressió, covariància, coeficient de correlació lineal, rectes de regressió}]{hyperref}

\definecolor{navy}{HTML}{1B2748}
\definecolor{accent}{HTML}{E8603A}
\definecolor{teal}{HTML}{1F9E90}
\definecolor{soft}{HTML}{EEF1F8}
\definecolor{muted}{HTML}{5D6685}

\setlength{\parindent}{0pt}
\setlength{\parskip}{3pt}
\settasks{label=\alph*), label-format=\bfseries\color{navy}, label-width=1.3em, item-indent=1.8em, column-sep=0.6em, after-item-skip=5pt, before-skip=3pt, after-skip=4pt}

% ---------- capçalera i peu ----------
\newcommand{\llicencia}{\href{https://creativecommons.org/licenses/by-nc-sa/4.0/deed.ca}{\ccbyncsa}\ \ Llicència Creative Commons Reconeixement-NoComercial-CompartirIgual 4.0 Internacional (CC BY-NC-SA 4.0)}
\pagestyle{fancy}
\fancyhf{}
\fancyhead[L]{\raisebox{-0.35\height}{\includegraphics[height=0.95cm]{logo-gva.png}}}
\fancyhead[C]{\small\color{muted}Solucions de la fitxa · Unitat 8: Distribucions bidimensionals · 1r Batxillerat CCSS}
\fancyhead[R]{\raisebox{-0.35\height}{\includegraphics[height=0.95cm]{logo-vermellar.png}}}
\fancyfoot[L]{\scriptsize\color{muted}\llicencia}
\fancyfoot[R]{\small\color{muted}\thepage}
\renewcommand{\headrulewidth}{0.4pt}
\renewcommand{\headrule}{\color{muted!40}\hrule height \headrulewidth}
\fancypagestyle{primera}{\fancyhf{}\renewcommand{\headrulewidth}{0pt}
  \fancyfoot[L]{\scriptsize\color{muted}\llicencia}
  \fancyfoot[R]{\small\color{muted}\thepage}}

% ---------- nivells i exercicis ----------
\newcommand{\estrelles}[1]{{\color{accent}\ifcase#1\or\faStar\faStar[regular]\faStar[regular]\or\faStar\faStar\faStar[regular]\or\faStar\faStar\faStar\fi}}
\newcounter{exer}
\newcommand{\ex}[1]{\par\needspace{4\baselineskip}\medskip\refstepcounter{exer}\noindent{\color{navy}\bfseries\theexer.}\ \ }

\newcommand{\seccio}[2]{\par\needspace{16\baselineskip}\bigskip
  \begin{tcolorbox}[colback=navy,colframe=navy,coltext=white,arc=2mm,boxrule=0pt,left=3mm,right=3mm,top=1.2mm,bottom=1.2mm]
  {\large\bfseries #1\quad #2}
  \end{tcolorbox}\vspace{-2pt}}

\newcommand{\exf}[1]{\par\medskip\refstepcounter{exer}\noindent{\color{navy}\bfseries\theexer.}\ \ }
\providecommand{\gt}{>}\providecommand{\lt}{<}
% quadre de recordatori (només a la fitxa d'exercicis)
\newcommand{\recorda}[1]{\par\smallskip\begin{tcolorbox}[colback=accent!7,colframe=accent!55,arc=2mm,boxrule=0.5pt,left=2.5mm,right=2.5mm,top=1mm,bottom=1mm]\small{\color{accent}\faLightbulb[regular]}\ \textbf{Recorda:}\ #1\end{tcolorbox}}
\newcommand{\res}[1]{\boxed{#1}}\newcommand{\rest}[1]{\fbox{#1}}

\def\fmneed{11}% que el títol d'una secció no quede sol al final del full
% =====================================================================
% Capçalera FentMates de les fitxes: un color per a cada curs.
% Els colors són els de la pàgina inicial (LEVELS a assets/js/materials.js);
% el logo ve de _eines/logo/pdf/ (python3 _eines/logo/genera.py).
%
% Ús, al final del preàmbul de la unitat:
%   % =====================================================================
% Capçalera FentMates de les fitxes: un color per a cada curs.
% Els colors són els de la pàgina inicial (LEVELS a assets/js/materials.js);
% el logo ve de _eines/logo/pdf/ (python3 _eines/logo/genera.py).
%
% Ús, al final del preàmbul de la unitat:
%   % =====================================================================
% Capçalera FentMates de les fitxes: un color per a cada curs.
% Els colors són els de la pàgina inicial (LEVELS a assets/js/materials.js);
% el logo ve de _eines/logo/pdf/ (python3 _eines/logo/genera.py).
%
% Ús, al final del preàmbul de la unitat:
%   \input{../comu/fentmates.tex}
%   \fmcurs{1eso}
% i, al començament del document:
%   \capcalera{Matemàtiques}{1r ESO}{Fitxa d'exercicis · Unitat 2}{Divisibilitat}
% Opcional abans de \input: \def\fmneed{14} (espai mínim abans de cada secció).
% =====================================================================
\makeatletter
\def\fm@logodir{../../logo/pdf}
\providecommand{\fmneed}{10}
\def\fm@cursnom{}

% clau del curs, color i color del text damunt del color
\newcommand{\fm@curs}[3]{\@namedef{fm@color@#1}{#2}\@namedef{fm@text@#1}{#3}}
\fm@curs{1eso}{E6B84A}{navy}
\fm@curs{2eso}{3B82F6}{white}
\fm@curs{3eso}{FF7A50}{white}
\fm@curs{4eso-a}{A78BFA}{white}
\fm@curs{4eso-b}{6D28D9}{white}
\fm@curs{1batx-mi}{9B1C1C}{white}
\fm@curs{1batx-ccss}{EF6B6B}{white}
\fm@curs{2batx-mii}{1E6B47}{white}
\fm@curs{2batx-ccss}{2BB3A3}{white}

\newcommand{\fmcurs}[1]{%
  \@ifundefined{fm@color@#1}{\PackageError{fentmates}{Curs desconegut: #1}{}}{}%
  \edef\fm@tmp{\noexpand\definecolor{curs}{HTML}{\@nameuse{fm@color@#1}}%
               \noexpand\colorlet{curstext}{\@nameuse{fm@text@#1}}}\fm@tmp
  \colorlet{cursfosc}{curs!72!black}%
  \def\fm@logo{\fm@logodir/vertical-#1.pdf}%
  \def\fm@marca{\fm@logodir/marca-#1.pdf}}

% ---------- capçalera de la primera pàgina ----------
% #1 matèria (petit), #2 curs, #3 tipus de document i unitat, #4 títol
\newcommand{\capcalera}[4]{%
  \gdef\fm@cursnom{#2}%
  \par\noindent
  \begin{tcolorbox}[colback=curs,colframe=curs,boxrule=0pt,arc=2.2mm,width=3.5cm,height=2.15cm,
      valign=center,left=3mm,right=2mm,top=1mm,bottom=1mm,nobeforeafter,box align=center]
    \color{curstext}{\footnotesize #1}\par\vspace{1pt}{\fontsize{19}{21}\selectfont\bfseries #2\par}%
  \end{tcolorbox}\hspace{5mm}%
  \begin{minipage}[c]{\dimexpr\linewidth-3.5cm-5mm-2.9cm\relax}
    {\color{cursfosc}\small\bfseries\MakeUppercase{#3}}\par\vspace{2pt}
    {\color{navy}\fontsize{20}{23}\selectfont\bfseries #4\par}
  \end{minipage}\hfill
  \begin{minipage}[c]{2.6cm}\raggedleft\includegraphics[width=2.6cm]{\fm@logo}\end{minipage}\par
  \vspace{2pt}}

% ---------- barres de secció del color del curs ----------
\renewcommand{\seccio}[2]{\par\needspace{\fmneed\baselineskip}\bigskip
  \begin{tcolorbox}[colback=curs,colframe=curs,coltext=curstext,arc=2mm,boxrule=0pt,left=3mm,right=3mm,top=1.2mm,bottom=1.2mm]
  {\large\bfseries #1\quad #2}
  \end{tcolorbox}\vspace{-2pt}}

% ---------- estrelles de dificultat del color del curs (a Batxillerat no se'n posen) ----------
\@ifundefined{estrelles}{}{%
  \renewcommand{\estrelles}[1]{{\color{curs}\ifcase#1\or\faStar\faStar[regular]\faStar[regular]\or\faStar\faStar\faStar[regular]\or\faStar\faStar\faStar\fi}}}

% ---------- capçalera i peu de les pàgines ----------
% primera pàgina: logos institucionals al peu
\fancypagestyle{primera}{\fancyhf{}\renewcommand{\headrulewidth}{0pt}
  \fancyfoot[L]{\raisebox{-0.32\height}{\includegraphics[height=0.78cm]{logo-gva.png}}\hspace{5mm}%
    \raisebox{-0.32\height}{\includegraphics[height=0.78cm]{logo-vermellar.png}}}
  \fancyfoot[R]{\scriptsize\color{muted}\href{https://creativecommons.org/licenses/by-nc-sa/4.0/deed.ca}{\ccbyncsa}\ \ CC BY-NC-SA 4.0 · fentmates.org\qquad\small\thepage}}
% la resta: el símbol a l'esquerra i el curs a la dreta, amb el títol al centre com abans
\fancyhead[L]{\raisebox{-0.3\height}{\includegraphics[height=0.62cm]{\fm@marca}}}
\fancyhead[R]{\small\bfseries\color{cursfosc}\fm@cursnom}
\renewcommand{\headrule}{\color{curs!45}\hrule height \headrulewidth}
\makeatother

%   \fmcurs{1eso}
% i, al començament del document:
%   \capcalera{Matemàtiques}{1r ESO}{Fitxa d'exercicis · Unitat 2}{Divisibilitat}
% Opcional abans de \input: \def\fmneed{14} (espai mínim abans de cada secció).
% =====================================================================
\makeatletter
\def\fm@logodir{../../logo/pdf}
\providecommand{\fmneed}{10}
\def\fm@cursnom{}

% clau del curs, color i color del text damunt del color
\newcommand{\fm@curs}[3]{\@namedef{fm@color@#1}{#2}\@namedef{fm@text@#1}{#3}}
\fm@curs{1eso}{E6B84A}{navy}
\fm@curs{2eso}{3B82F6}{white}
\fm@curs{3eso}{FF7A50}{white}
\fm@curs{4eso-a}{A78BFA}{white}
\fm@curs{4eso-b}{6D28D9}{white}
\fm@curs{1batx-mi}{9B1C1C}{white}
\fm@curs{1batx-ccss}{EF6B6B}{white}
\fm@curs{2batx-mii}{1E6B47}{white}
\fm@curs{2batx-ccss}{2BB3A3}{white}

\newcommand{\fmcurs}[1]{%
  \@ifundefined{fm@color@#1}{\PackageError{fentmates}{Curs desconegut: #1}{}}{}%
  \edef\fm@tmp{\noexpand\definecolor{curs}{HTML}{\@nameuse{fm@color@#1}}%
               \noexpand\colorlet{curstext}{\@nameuse{fm@text@#1}}}\fm@tmp
  \colorlet{cursfosc}{curs!72!black}%
  \def\fm@logo{\fm@logodir/vertical-#1.pdf}%
  \def\fm@marca{\fm@logodir/marca-#1.pdf}}

% ---------- capçalera de la primera pàgina ----------
% #1 matèria (petit), #2 curs, #3 tipus de document i unitat, #4 títol
\newcommand{\capcalera}[4]{%
  \gdef\fm@cursnom{#2}%
  \par\noindent
  \begin{tcolorbox}[colback=curs,colframe=curs,boxrule=0pt,arc=2.2mm,width=3.5cm,height=2.15cm,
      valign=center,left=3mm,right=2mm,top=1mm,bottom=1mm,nobeforeafter,box align=center]
    \color{curstext}{\footnotesize #1}\par\vspace{1pt}{\fontsize{19}{21}\selectfont\bfseries #2\par}%
  \end{tcolorbox}\hspace{5mm}%
  \begin{minipage}[c]{\dimexpr\linewidth-3.5cm-5mm-2.9cm\relax}
    {\color{cursfosc}\small\bfseries\MakeUppercase{#3}}\par\vspace{2pt}
    {\color{navy}\fontsize{20}{23}\selectfont\bfseries #4\par}
  \end{minipage}\hfill
  \begin{minipage}[c]{2.6cm}\raggedleft\includegraphics[width=2.6cm]{\fm@logo}\end{minipage}\par
  \vspace{2pt}}

% ---------- barres de secció del color del curs ----------
\renewcommand{\seccio}[2]{\par\needspace{\fmneed\baselineskip}\bigskip
  \begin{tcolorbox}[colback=curs,colframe=curs,coltext=curstext,arc=2mm,boxrule=0pt,left=3mm,right=3mm,top=1.2mm,bottom=1.2mm]
  {\large\bfseries #1\quad #2}
  \end{tcolorbox}\vspace{-2pt}}

% ---------- estrelles de dificultat del color del curs (a Batxillerat no se'n posen) ----------
\@ifundefined{estrelles}{}{%
  \renewcommand{\estrelles}[1]{{\color{curs}\ifcase#1\or\faStar\faStar[regular]\faStar[regular]\or\faStar\faStar\faStar[regular]\or\faStar\faStar\faStar\fi}}}

% ---------- capçalera i peu de les pàgines ----------
% primera pàgina: logos institucionals al peu
\fancypagestyle{primera}{\fancyhf{}\renewcommand{\headrulewidth}{0pt}
  \fancyfoot[L]{\raisebox{-0.32\height}{\includegraphics[height=0.78cm]{logo-gva.png}}\hspace{5mm}%
    \raisebox{-0.32\height}{\includegraphics[height=0.78cm]{logo-vermellar.png}}}
  \fancyfoot[R]{\scriptsize\color{muted}\href{https://creativecommons.org/licenses/by-nc-sa/4.0/deed.ca}{\ccbyncsa}\ \ CC BY-NC-SA 4.0 · fentmates.org\qquad\small\thepage}}
% la resta: el símbol a l'esquerra i el curs a la dreta, amb el títol al centre com abans
\fancyhead[L]{\raisebox{-0.3\height}{\includegraphics[height=0.62cm]{\fm@marca}}}
\fancyhead[R]{\small\bfseries\color{cursfosc}\fm@cursnom}
\renewcommand{\headrule}{\color{curs!45}\hrule height \headrulewidth}
\makeatother

%   \fmcurs{1eso}
% i, al començament del document:
%   \capcalera{Matemàtiques}{1r ESO}{Fitxa d'exercicis · Unitat 2}{Divisibilitat}
% Opcional abans de \input: \def\fmneed{14} (espai mínim abans de cada secció).
% =====================================================================
\makeatletter
\def\fm@logodir{../../logo/pdf}
\providecommand{\fmneed}{10}
\def\fm@cursnom{}

% clau del curs, color i color del text damunt del color
\newcommand{\fm@curs}[3]{\@namedef{fm@color@#1}{#2}\@namedef{fm@text@#1}{#3}}
\fm@curs{1eso}{E6B84A}{navy}
\fm@curs{2eso}{3B82F6}{white}
\fm@curs{3eso}{FF7A50}{white}
\fm@curs{4eso-a}{A78BFA}{white}
\fm@curs{4eso-b}{6D28D9}{white}
\fm@curs{1batx-mi}{9B1C1C}{white}
\fm@curs{1batx-ccss}{EF6B6B}{white}
\fm@curs{2batx-mii}{1E6B47}{white}
\fm@curs{2batx-ccss}{2BB3A3}{white}

\newcommand{\fmcurs}[1]{%
  \@ifundefined{fm@color@#1}{\PackageError{fentmates}{Curs desconegut: #1}{}}{}%
  \edef\fm@tmp{\noexpand\definecolor{curs}{HTML}{\@nameuse{fm@color@#1}}%
               \noexpand\colorlet{curstext}{\@nameuse{fm@text@#1}}}\fm@tmp
  \colorlet{cursfosc}{curs!72!black}%
  \def\fm@logo{\fm@logodir/vertical-#1.pdf}%
  \def\fm@marca{\fm@logodir/marca-#1.pdf}}

% ---------- capçalera de la primera pàgina ----------
% #1 matèria (petit), #2 curs, #3 tipus de document i unitat, #4 títol
\newcommand{\capcalera}[4]{%
  \gdef\fm@cursnom{#2}%
  \par\noindent
  \begin{tcolorbox}[colback=curs,colframe=curs,boxrule=0pt,arc=2.2mm,width=3.5cm,height=2.15cm,
      valign=center,left=3mm,right=2mm,top=1mm,bottom=1mm,nobeforeafter,box align=center]
    \color{curstext}{\footnotesize #1}\par\vspace{1pt}{\fontsize{19}{21}\selectfont\bfseries #2\par}%
  \end{tcolorbox}\hspace{5mm}%
  \begin{minipage}[c]{\dimexpr\linewidth-3.5cm-5mm-2.9cm\relax}
    {\color{cursfosc}\small\bfseries\MakeUppercase{#3}}\par\vspace{2pt}
    {\color{navy}\fontsize{20}{23}\selectfont\bfseries #4\par}
  \end{minipage}\hfill
  \begin{minipage}[c]{2.6cm}\raggedleft\includegraphics[width=2.6cm]{\fm@logo}\end{minipage}\par
  \vspace{2pt}}

% ---------- barres de secció del color del curs ----------
\renewcommand{\seccio}[2]{\par\needspace{\fmneed\baselineskip}\bigskip
  \begin{tcolorbox}[colback=curs,colframe=curs,coltext=curstext,arc=2mm,boxrule=0pt,left=3mm,right=3mm,top=1.2mm,bottom=1.2mm]
  {\large\bfseries #1\quad #2}
  \end{tcolorbox}\vspace{-2pt}}

% ---------- estrelles de dificultat del color del curs (a Batxillerat no se'n posen) ----------
\@ifundefined{estrelles}{}{%
  \renewcommand{\estrelles}[1]{{\color{curs}\ifcase#1\or\faStar\faStar[regular]\faStar[regular]\or\faStar\faStar\faStar[regular]\or\faStar\faStar\faStar\fi}}}

% ---------- capçalera i peu de les pàgines ----------
% primera pàgina: logos institucionals al peu
\fancypagestyle{primera}{\fancyhf{}\renewcommand{\headrulewidth}{0pt}
  \fancyfoot[L]{\raisebox{-0.32\height}{\includegraphics[height=0.78cm]{logo-gva.png}}\hspace{5mm}%
    \raisebox{-0.32\height}{\includegraphics[height=0.78cm]{logo-vermellar.png}}}
  \fancyfoot[R]{\scriptsize\color{muted}\href{https://creativecommons.org/licenses/by-nc-sa/4.0/deed.ca}{\ccbyncsa}\ \ CC BY-NC-SA 4.0 · fentmates.org\qquad\small\thepage}}
% la resta: el símbol a l'esquerra i el curs a la dreta, amb el títol al centre com abans
\fancyhead[L]{\raisebox{-0.3\height}{\includegraphics[height=0.62cm]{\fm@marca}}}
\fancyhead[R]{\small\bfseries\color{cursfosc}\fm@cursnom}
\renewcommand{\headrule}{\color{curs!45}\hrule height \headrulewidth}
\makeatother

\fmcurs{1batx-ccss}

\begin{document}
\thispagestyle{primera}

\capcalera{Matemàtiques CCSS}{1r Batx}{Solucions de la fitxa · Unitat 8}{\raggedright Distribucions bidimensionals}
\vspace{8pt}

\begin{tcolorbox}[colback=soft,colframe=soft,arc=2mm,left=3mm,right=3mm,top=1.5mm,bottom=1.5mm]
\small Intenta resoldre cada exercici abans de mirar-ne la solució. Si no t'ix el mateix resultat, repassa els passos o comprova'l amb les calculadores dels apunts interactius de la unitat.
\end{tcolorbox}

% =====================================================
\seccio{1}{Distribucions bidimensionals. Taules de freqüències}

\exf{1} \par{\centering {\small$\begin{array}{|c|c|c|c|c|}\hline X\backslash Y & 0 & 1 & 2 & f_x\\\hline 1 & 4 & 2 & 0 & 6\\ 2 & 3 & 6 & 3 & 12\\ 3 & 0 & 3 & 4 & 7\\\hline f_y & 7 & 11 & 7 & 25\\\hline \end{array}$}\par}\smallskip\begin{tasks}(1)
  \task Vegeu la taula.
  \task $\res{6}$ famílies amb 2 fills i 1 cotxe; \ $\res{7}$ amb més d'un cotxe (2 cotxes).
  \task $\bar x=\frac{51}{25}=\res{2{,}04}$, \ $\bar y=\frac{25}{25}=\res{1}$
\end{tasks}

\ex{2} \par{\centering {\small$\begin{array}{|c|c|c|c|c|}\hline X\backslash Y & 1 & 2 & 3 & f_x\\\hline 1 & 2 & 1 & 0 & 3\\ 2 & 2 & 3 & 0 & 5\\ 3 & 0 & 2 & 2 & 4\\\hline f_y & 4 & 6 & 2 & 12\\\hline \end{array}$}\par}\smallskip $\bar x=\frac{25}{12}\approx\res{2{,}08}$, \ $\bar y=\frac{22}{12}\approx\res{1{,}83}$.

\ex{2} \par{\centering {\small$\begin{array}{|c|c|c|c|c|}\hline X\backslash Y & {[0,5)} & {[5,7)} & {[7,10]} & f_x\\\hline {[0,2)} & 8 & 4 & 1 & 13\\ {[2,4)} & 5 & 10 & 6 & 21\\ {[4,6)} & 1 & 5 & 10 & 16\\\hline f_y & 14 & 19 & 17 & 50\\\hline \end{array}$}\par}\smallskip\begin{tasks}(1)
  \task Marques de classe: $X$: $1,\ 3,\ 5$; \ $Y$: $2{,}5;\ 6;\ 8{,}5$. \ $\bar x=\frac{156}{50}=\res{3{,}12}$ h, \ $\bar y=\frac{293{,}5}{50}=\res{5{,}87}$
  \task $\frac{6}{50}=0{,}12$: \ $\res{12\,\%}$
  \task $\frac{10}{16}=\res{62{,}5\,\%}$; \ dels que estudien menys de 2 h, $\frac{1}{13}\approx\res{7{,}7\,\%}$.
  \task $\frac{8}{14}\approx\res{57{,}1\,\%}$
\end{tasks}


% =====================================================
\seccio{2}{Núvol de punts}

\exf{1} $(\bar x,\bar y)=(4{,}5;\ 51{,}25)$. \ Els punts baixen gairebé en línia recta: correlació \rest{negativa i molt forta}.\par\smallskip{\centering \begin{tikzpicture}[scale=0.42,baseline=(current bounding box.center)]\draw[muted!30,very thin] (0,0) grid (9.0,6.0);\draw[-latex,thick] (0,0)--(9.5,0) node[right,font=\scriptsize]{$h$};\draw[-latex,thick] (0,0)--(0,6.5) node[above,font=\scriptsize]{$l$};\node[below,font=\tiny] at (0.0,0) {$0$};\node[below,font=\tiny] at (1.0,0) {$1$};\node[below,font=\tiny] at (2.0,0) {$2$};\node[below,font=\tiny] at (3.0,0) {$3$};\node[below,font=\tiny] at (4.0,0) {$4$};\node[below,font=\tiny] at (5.0,0) {$5$};\node[below,font=\tiny] at (6.0,0) {$6$};\node[below,font=\tiny] at (7.0,0) {$7$};\node[below,font=\tiny] at (8.0,0) {$8$};\node[below,font=\tiny] at (9.0,0) {$9$};\node[left,font=\tiny] at (0,0.0) {$35$};\node[left,font=\tiny] at (0,1.0) {$40$};\node[left,font=\tiny] at (0,2.0) {$45$};\node[left,font=\tiny] at (0,3.0) {$50$};\node[left,font=\tiny] at (0,4.0) {$55$};\node[left,font=\tiny] at (0,5.0) {$60$};\node[left,font=\tiny] at (0,6.0) {$65$};\fill[accent] (1.0,5.4) circle (0.16);\fill[accent] (2.0,5.0) circle (0.16);\fill[accent] (3.0,4.0) circle (0.16);\fill[accent] (4.0,3.6) circle (0.16);\fill[accent] (5.0,3.0) circle (0.16);\fill[accent] (6.0,2.2) circle (0.16);\fill[accent] (7.0,1.8) circle (0.16);\fill[accent] (8.0,1.0) circle (0.16);\draw[navy,thick,fill=white] (9/2,13/4) circle (0.2);\end{tikzpicture}\par}

\ex{1} \begin{tasks}(4)
  \task $\res{r=0{,}9}$
  \task $\res{r=-0{,}95}$
  \task $\res{r=0}$
  \task $\res{r=-0{,}5}$
\end{tasks}

\ex{1} \begin{tasks}(2)
  \task \rest{positiva}
  \task \rest{negativa}
  \task \rest{cap}
  \task \rest{positiva}
  \task \rest{negativa}
  \task \rest{positiva}
\end{tasks}


% =====================================================
\seccio{3}{Correlació i regressió. Tipus de dependència}

\exf{1} \begin{tasks}(3)
  \task \rest{funcional} ($y=4x$)
  \task \rest{estadística}
  \task \rest{independents}
  \task \rest{estadística}
  \task \rest{funcional}
  \task \rest{estadística}
\end{tasks}

\ex{1} \begin{tasks}(1)
  \task \rest{Funcional}: els punts són exactament sobre una corba (no és una recta: la regressió lineal no és l'adequada).
  \task \rest{Estadística}: núvol allargat al voltant d'una recta creixent. Ací sí que té sentit buscar una recta de regressió.
  \task \rest{Independència}: el núvol no té cap forma.
\end{tasks}

\ex{2} \rest{No}: la correlació no implica causalitat. Hi ha una tercera variable, la \rest{grandària de l'incendi}: com més gran és, més bombers hi van i més danys causa.


% =====================================================
\seccio{4}{Correlació lineal. Covariància}

\exf{2} \par{\centering {\small$\begin{array}{|c|c|c|c|c|}\hline x_i & y_i & x_i^2 & y_i^2 & x_iy_i\\\hline 2 & 3 & 4 & 9 & 6\\ 4 & 5 & 16 & 25 & 20\\ 5 & 4 & 25 & 16 & 20\\ 6 & 7 & 36 & 49 & 42\\ 8 & 9 & 64 & 81 & 72\\\hline \textstyle\sum: 25 & 28 & 145 & 180 & 160\\\hline\end{array}$}\par}\smallskip $\bar x=\frac{25}{5}=5$, \ $\bar y=\frac{28}{5}=5{,}6$; \ $\sigma_x^2=\frac{145}{5}-\bar x^2=4$, \ $\sigma_y^2=\frac{180}{5}-\bar y^2=4{,}64$; \ $\sigma_x=2$, \ $\sigma_y\approx2{,}154$. \ $\sigma_{xy}=\frac{160}{5}-\bar x\,\bar y=\res{4}$.

\ex{2} $\bar x=\frac{150}{5}=30$, \ $\bar y=\frac{30}{5}=6$; \ $\sigma_x^2=\frac{5500}{5}-\bar x^2=200$, \ $\sigma_y^2=\frac{210}{5}-\bar y^2=6$; \ $\sigma_x\approx14{,}142$, \ $\sigma_y\approx2{,}449$. \ $\sigma_{xy}=\frac{730}{5}-\bar x\,\bar y=\res{-34}$. \ És negativa: quan $X$ augmenta, $Y$ tendeix a \rest{disminuir}.

\ex{2} $\sum x_iy_jf_{ij}=50$: \ $\sigma_{xy}=\frac{50}{12}-\bar x\,\bar y\approx\res{0{,}347}$

\ex{2} $\sum x_iy_jf_{ij}=59$: \ $\sigma_{xy}=\frac{59}{25}-2{,}04\cdot1=\res{0{,}32}$

\ex{3} $\sum x_i^2f_x=602$, \ $\sum y_j^2f_y=1999{,}75$, \ $\sum x_iy_jf_{ij}=1010{,}5$. \ $\sigma_x^2=\frac{602}{50}-3{,}12^2\approx\res{2{,}306}$, \ $\sigma_y^2=\frac{1999{,}75}{50}-5{,}87^2\approx\res{5{,}538}$, \ $\sigma_{xy}=\frac{1010{,}5}{50}-3{,}12\cdot5{,}87\approx\res{1{,}896}$. \ És positiva: \rest{com més hores d'estudi, més nota} (tendència).


% =====================================================
\seccio{5}{Coeficient de correlació lineal}

\exf{2} \begin{tasks}(1)
  \task $r=\dfrac{\sigma_{xy}}{\sigma_x\sigma_y}\approx\res{0{,}928}$: correlació \rest{positiva forta}.
  \task $r=\dfrac{\sigma_{xy}}{\sigma_x\sigma_y}\approx\res{-0{,}981}$: correlació \rest{negativa molt forta}.
  \task $r=\dfrac{0{,}32}{\sqrt{0{,}518}\cdot\sqrt{0{,}56}}\approx\res{0{,}594}$: correlació \rest{positiva moderada}.
  \task $r=\dfrac{1{,}896}{\sqrt{2{,}306}\cdot\sqrt{5{,}538}}\approx\res{0{,}53}$: correlació \rest{positiva moderada}.
\end{tasks}

\ex{2} $\bar x=\frac{1155}{7}=165$, \ $\bar y=\frac{421}{7}\approx60{,}143$; \ $\sigma_x^2=\frac{191\,275}{7}-\bar x^2=100$, \ $\sigma_y^2=\frac{25\,807}{7}-\bar y^2\approx69{,}551$; \ $\sigma_x=10$, \ $\sigma_y\approx8{,}34$. \ $\sigma_{xy}=\frac{70\,045}{7}-\bar x\,\bar y\approx\res{82{,}857}$. $r=\dfrac{\sigma_{xy}}{\sigma_x\sigma_y}\approx\res{0{,}994}$: correlació \rest{positiva molt forta}.\par\smallskip{\centering \begin{tikzpicture}[scale=0.42,baseline=(current bounding box.center)]\draw[muted!30,very thin] (0,0) grid (8.0,6.0);\draw[-latex,thick] (0,0)--(8.5,0) node[right,font=\scriptsize]{$cm$};\draw[-latex,thick] (0,0)--(0,6.5) node[above,font=\scriptsize]{$kg$};\node[below,font=\tiny] at (0.0,0) {$145$};\node[below,font=\tiny] at (2.0,0) {$155$};\node[below,font=\tiny] at (4.0,0) {$165$};\node[below,font=\tiny] at (6.0,0) {$175$};\node[below,font=\tiny] at (8.0,0) {$185$};\node[left,font=\tiny] at (0,0.0) {$45$};\node[left,font=\tiny] at (0,1.0) {$50$};\node[left,font=\tiny] at (0,2.0) {$55$};\node[left,font=\tiny] at (0,3.0) {$60$};\node[left,font=\tiny] at (0,4.0) {$65$};\node[left,font=\tiny] at (0,5.0) {$70$};\node[left,font=\tiny] at (0,6.0) {$75$};\fill[accent] (1.0,0.6) circle (0.16);\fill[accent] (2.0,1.4) circle (0.16);\fill[accent] (3.0,2.0) circle (0.16);\fill[accent] (4.0,3.2) circle (0.16);\fill[accent] (5.0,3.6) circle (0.16);\fill[accent] (6.0,5.0) circle (0.16);\fill[accent] (7.0,5.4) circle (0.16);\end{tikzpicture}\par}

\ex{1} Tots els $x_i$ es divideixen entre 100 i la covariància també: $\sigma_{xy}\approx\res{0{,}8286}$. En canvi, $r$ \rest{no canvia} ($r\approx0{,}994$), perquè no depén de les unitats.

\ex{2} $\bar x=\frac{15}{6}=2{,}5$, \ $\bar y=\frac{40}{6}\approx6{,}667$; \ $\sigma_x^2=\frac{55}{6}-\bar x^2\approx2{,}917$, \ $\sigma_y^2=\frac{306}{6}-\bar y^2\approx6{,}556$; \ $\sigma_x\approx1{,}708$, \ $\sigma_y\approx2{,}56$. \ $\sigma_{xy}=\frac{75}{6}-\bar x\,\bar y\approx\res{-4{,}167}$. $r=\dfrac{\sigma_{xy}}{\sigma_x\sigma_y}\approx\res{-0{,}953}$: correlació \rest{negativa molt forta}. \ Però una correlació, encara que siga forta, \rest{no demostra que hi haja una relació de causa i efecte}: hi pot haver altres variables (el temps d'estudi, per exemple).

\ex{2} $\bar x=\frac{153}{6}=25{,}5$, \ $\bar y=\frac{406}{6}\approx67{,}667$; \ $\sigma_x^2=\frac{4059}{6}-\bar x^2=26{,}25$, \ $\sigma_y^2=\frac{30\,378}{6}-\bar y^2\approx484{,}222$; \ $\sigma_x\approx5{,}123$, \ $\sigma_y\approx22{,}005$. \ $\sigma_{xy}=\frac{11\,025}{6}-\bar x\,\bar y=\res{112}$. $r=\dfrac{\sigma_{xy}}{\sigma_x\sigma_y}\approx\res{0{,}993}$: correlació \rest{positiva molt forta}.

\ex{2} $\sigma_x^2\approx0{,}576$, \ $\sigma_y^2\approx0{,}472$, \ $\sigma_{xy}\approx0{,}347$. \ $r\approx\res{0{,}666}$: correlació \rest{positiva moderada}.

\ex{1} $\sigma_{xy}=r\,\sigma_x\,\sigma_y=0{,}8\cdot2\cdot5=\res{8}$.

\ex{1} $r=\frac{-12}{4\cdot5}=\res{-0{,}6}$ (negativa moderada). \ En el segon cas eixiria $r=\frac{3}{2}=1{,}5>1$: \rest{no és possible}, hi ha un error en les dades.

\ex{3} $\bar x=\frac{50}{10}=5$, \ $\bar y=\frac{80}{10}=8$; \ $\sigma_x^2=\frac{290}{10}-\bar x^2=4$, \ $\sigma_y^2=\frac{730}{10}-\bar y^2=9$; \ $\sigma_x=2$, \ $\sigma_y=3$. \ $\sigma_{xy}=\frac{440}{10}-\bar x\,\bar y=\res{4}$. $r=\dfrac{\sigma_{xy}}{\sigma_x\sigma_y}\approx\res{0{,}667}$: correlació \rest{positiva moderada}.

\ex{2} \begin{tasks}(1)
  \task Relació \rest{parabòlica} (creix i després decreix): és forta, però no és lineal i $r\approx0{,}02$, \rest{pròxim a 0}.
  \task Relació \rest{exponencial}: creixent, $r\approx0{,}9$ és alt, però els punts no segueixen una recta.
  \task Els punts de baix no tenen cap relació, però el \rest{valor atípic} $(9,9)$ fa que $r\approx0{,}99$. Cal revisar-lo abans de traure conclusions.
\end{tasks}

\ex{1} a) \rest{Falsa}: només vol dir que no hi ha relació lineal (pot ser parabòlica, per exemple). \ b) \rest{Falsa}: és molt forta; el signe indica que és inversa. \ c) \rest{Vertadera}.


% =====================================================
\seccio{6}{Rectes de regressió}

\exf{2} $Y$ sobre $X$: $y-5{,}6=\frac{4}{4}(x-5)\Rightarrow\res{y=x+0{,}6}$; \ $\hat y(7)=\res{7{,}6}$. \ $X$ sobre $Y$: $x-5=\frac{4}{4{,}64}(y-5{,}6)\Rightarrow\res{x\approx0{,}862y+0{,}172}$; \ $\hat x(6)\approx\res{5{,}345}$. \ Són fiables: $r\approx0{,}928$ i els valors són dins de les dades.

\ex{2} \begin{tasks}(1)
  \task $\bar x=\frac{36}{8}=4{,}5$, \ $\bar y=\frac{158}{8}=19{,}75$; \ $\sigma_x^2=\frac{204}{8}-\bar x^2=5{,}25$, \ $\sigma_y^2=\frac{3328}{8}-\bar y^2\approx25{,}938$; \ $\sigma_x\approx2{,}291$, \ $\sigma_y\approx5{,}093$. \ $\sigma_{xy}=\frac{803}{8}-\bar x\,\bar y=\res{11{,}5}$. $r=\dfrac{\sigma_{xy}}{\sigma_x\sigma_y}\approx\res{0{,}985}$: correlació \rest{positiva molt forta}.
  \task $y-19{,}75=\frac{11{,}5}{5{,}25}(x-4{,}5)\Rightarrow\res{y\approx2{,}19x+9{,}893}$. Per cada 1000 € més de publicitat, les vendes augmenten, de mitjana, uns $\res{2190}$ €.
  \task $\hat y(5{,}5)\approx21{,}94$: unes vendes de $\res{21\,940}$ €.
  \task Recta de $X$ sobre $Y$: $x-4{,}5=\frac{11{,}5}{25{,}938}(y-19{,}75)\Rightarrow\res{x\approx0{,}443y-4{,}257}$; \ $\hat x(24)\approx6{,}384$: cal invertir uns $\res{6384}$ €. Fiable: $r\approx0{,}985$.
\end{tasks}

\ex{2} $y-60{,}143=\frac{82{,}857}{100}(x-165)\Rightarrow\res{y\approx0{,}829x-76{,}571}$. \ $\hat y(168)\approx\res{62{,}63\text{ kg}}$: \rest{fiable} ($r\approx0{,}994$ i 168 és dins de les dades). \ $\hat y(120)\approx22{,}86$ kg: \rest{no és fiable}, 120 cm queda molt lluny de l'interval $[150,180]$. \ Altura sobre pes: $\res{x\approx1{,}191y+93{,}351}$; \ $\hat x(65)\approx\res{170{,}8\text{ cm}}$.

\ex{2} Recta: $\res{y\approx4{,}267x-41{,}133}$. \ $\hat y(25)\approx65{,}53$: uns $\res{66}$ gelats. Molt fiable: $r\approx0{,}993$.

\ex{2} $y-1=\frac{0{,}32}{0{,}518}(x-2{,}04)\Rightarrow\res{y\approx0{,}617x-0{,}259}$; \ $\hat y(2)\approx\res{0{,}975}$ cotxes. Poc fiable: $r\approx0{,}594$ només indica una correlació moderada.

\ex{2} $\bar x=5$, $\bar y=8$, $\sigma_x^2=4$, $\sigma_{xy}=4$: \ $y-8=\frac44(x-5)\Rightarrow\res{y=x+3}$; \ $\hat y(6)=\res{9}$ (poc precisa: $r\approx0{,}667$ és moderat).

\ex{2} $\sigma_{xy}=r\,\sigma_x\sigma_y=-0{,}75\cdot2\cdot4=-6$. \ Y sobre X: $y-15=\frac{-6}{4}(x-8)\Rightarrow\res{y=-1{,}5x+27}$; \ $\hat y(9)=\res{13{,}5}$. \ X sobre Y: $x-8=\frac{-6}{16}(y-15)\Rightarrow\res{x=-0{,}375y+13{,}625}$; \ $\hat x(12)=\res{9{,}125}$.

\ex{3} Substituint: $x=1{,}2(0{,}5x+2)-1=0{,}6x+1{,}4\Rightarrow\bar x=\res{3{,}5}$, \ $\bar y=\res{3{,}75}$. \ $r^2=0{,}5\cdot1{,}2=0{,}6$ i, com que els coeficients són positius, $r=\sqrt{0{,}6}\approx\res{0{,}775}$.

\ex{3} Resolent el sistema: $(\bar x,\bar y)=\res{(5;\ 4)}$. \ En la forma $y=ax+b$: $y=0{,}8x$ i $y=2x-6$. La més plana, $\res{y=0{,}8x}$, és la de $Y$ sobre $X$ ($m_{yx}=0{,}8$); l'altra, $\res{x=0{,}5y+3}$, és la de $X$ sobre $Y$ ($m_{xy}=0{,}5$). (A l'inrevés eixiria $r^2=2\cdot1{,}25=2{,}5>1$.) \ $r^2=0{,}8\cdot0{,}5=0{,}4\Rightarrow r=\sqrt{0{,}4}\approx\res{0{,}632}$. \ $\hat y(6)=\res{4{,}8}$, \ $\hat x(5)=\res{5{,}5}$.

\ex{2} a) \rest{No}: els dos coeficients de regressió han de tindre el mateix signe (el de $\sigma_{xy}$), i ací un és positiu i l'altre negatiu. \ b) \rest{Sí}: tots dos són negatius i $r^2=(-0{,}4)(-0{,}9)=0{,}36\leq1$, \ $r=\res{-0{,}6}$.


\end{document}
